\section{知识记忆}\label{sec:knowledge}

\subsection{写入与检索知识}

知识记忆以「主语—谓语—宾语」三元组结构存储事实，助手通过 \code{kylin\_memory} 工具
读写。日常使用完全通过自然语言完成。要写入知识，直接说「记住：张三的工作邮箱是
zs@example.com」，入库前邮箱会被自动脱敏；要检索，就问「张三的邮箱是什么」或「查一下
和部署相关的记忆」；想了解某条事实的来龙去脉，可以问「这条事实是什么时候改成现在
这样的」，助手会展示完整的版本链。

检索时系统同时匹配字段、关键词与知识条目之间的关联关系，并对中文同义表述做归一化
（「住址」与「地址」视为同一谓语），因此不必刻意使用与当初写入时完全相同的措辞。
需要语义级的模糊检索（例如「找之前聊过的类似问题」）时，助手会改用语义检索工具
（第~\ref{sec:config}~章配置的嵌入后端）。

\subsection{新旧知识冲突时会发生什么}

当你写入的新事实与已有事实矛盾时，系统按明确的规则处置，而不是简单地覆盖或报错。
同一条事实的重复确认会增强其置信度；信息互含时会合并保留更完整的表述；住址、电话、
邮箱、工作单位、状态等天然互斥的谓词发生变更时，旧值自动转为\textbf{历史版本}并记录
失效时间，新值接替当前值；无法自动判定的矛盾则两者并存，等待后续信息澄清。全部
处置过程可追溯——每条事实都有完整的版本链，任何时刻的值都可以回放。

\subsection{把成功流程沉淀为模板}

当某类任务被成功完成过多次（例如同一类构建错误的修复流程），你可以让助手「把这个
流程保存为模板」。模板不是立即生效的：它先进入草稿状态，之后每次按此流程执行都会
积累一条真实执行证据；当证据满五条、成功率达到八成，并经你人工确认后，模板才被
激活供日后复用。这套门槛的目的是防止把偶然成功的做法固化成套路。激活后的模板如果
所依据的知识被删除（例如执行了遗忘），会自动退役，避免引用失效的经验。
